% !TeX spellcheck = en_US
%!TEX root=../Main.tex
% \clearpage
\section{Related Work}

\textbf{Image Foundation Models.}
Most of the current vision models are only suitable for specific tasks and domains,
and they require manually labeled datasets for training.
Regarding this, recent works have proposed vision foundation models. CLIP \citep{CLIP} and ALIGN \citep{align} prepare web-scale noisy image-text pairs to train dual-encoder models with contrastive learning,
leading to robust image-text representations for powerful zero-shot transfer. 
INTERN~\citep{shao2021intern} expands the self-supervised pretraining into multiple learning stages, which use a large quantity of image-text pairs as well as manually annotated images. INTERN achieves a better linear probe performance compared with CLIP, and improves data efficiency in the downstream image tasks.
Florence \citep{florence} extends them with unified contrastive learning \citep{unicl} and elaborate adaptation models, which support a wide range of vision tasks in different transfer settings. SimVLM \citep{simvlm} and OFA \citep{ofa} train encoder-decoder models with generative targets and show competitive performances on a series of multimodal tasks. Besides, CoCa \citep{coca} unifies contrastive learning as CLIP and generative learning as SimVLM. Recently, BeiT-3 \citep{beit3} introduces Multiway Transformers with unified BeiT \citep{beit} pretraining, achieving state-of-the-art transfer results on several vision and image-language tasks.

\textbf{Video Foundation Models.}
Previous image foundation models \citep{florence,coca} only show promising performance for video recognition (especially on Kinetics). As for video multimodal tasks, VIOLET \citep{violet} combines masked language and masked video modeling, All-in-one \citep{all_in_one} proposes unified video-language pretraining with a shared backbone, and LAVENDER \citep{lavender} unifies the tasks as masked language modeling. Though they perform well in multimodal benchmarks, they are trained with limited video-text data and struggle for video-only tasks, \eg\ action recognition. In contrast,
MERLOT Reserve \citep{reserve} collects 20M video-text-audio pairs to train the joint video representations with contrastive span matching, thus setting state-of-the-art video recognition and visual commonsense reasoning. Compared with image foundation models,
current video foundation models support limited video and video-language tasks,
especially for those fine-grained temporal discrimination tasks such as temporal localization.

\textbf{Self-supervised Pretraining.}
Self-supervised learning has developed rapidly recently. It focuses on designing different pretext tasks for pretraining \citep{doersch2015unsupervised,wang2015unsupervised,noroozi2016unsupervised,zhang2016colorful,feichtenhofer2022masked}, which can be mainly divided into contrastive learning and masked modeling. Contrastive learning adopts various data augmentations to generate different views of an image, 
then pulls together the positive pairs and pushes apart the negative pairs.
To maintain enough informative negative samples, previous methods depend on large memory banks or batch size \citep{wu2018unsupervised,he2020momentum,chen2020simple}.
BYOL \citep{byol} and SimSiam \citep{simsiam} eliminate the requirement of negative samples, designing elaborate techniques to avoid model collapse. As for masked modeling,
it learns rich visual representation via masked prediction based on visible context.
iGPT \citep{igpt} firstly mentions Masked Image Modeling (MIM). BeiT \citep{beit} propose visual token prediction with the pretrained tokenizer \citep{ramesh2021zero},
MaskFeat \citep{maskfeat} predicts the hand-crafted image descriptor, and MAE \citep{mae} directly reconstructs the raw pixels. For spatiotemporal representation learning,
VideoMAE \citep{videomae} and BEVT \citep{bevt} respectively extend MAE and BeiT to spatiotemporal space.

\textbf{Multimodal Pretraining.}
Starting from the development of image-text pretraining, large-scale video-text pretraining with specific downstream task finetuning has become the standard paradigm in the video-language area \citep{cpd,MiechASLSZ20,reserve,violet,videoclip,hu2022scaling,dou2022empirical,shen2021much,yao2021filip}. The seminal methods \citep{videobert,actbert} use pretrained visual and language encoders to extract the offline video and text features, while the recent methods \citep{cpd,MiechASLSZ20,clipbert,frozen,merlot,all_in_one} have demonstrated the feasibility of end-to-end training.  Besides, the popular methods often include two or three pretraining tasks,
\eg\, masked language modeling \citep{lavender},  video-text matching \citep{all_in_one}, video-text contrastive learning \citep{videoclip} and video-text masked modeling \citep{violet}.
